%
% 51-integral.tex -- Das Argumentintegral für f(z) = z^n
%
% (c) 2026 Prof Dr Andreas Müller
%
\subsection{Auswertung des Argumentintegrals
\label{buch:analytisch:argument:subsection:zn}}
In diesem Abschnitt soll die Verbindung zwischen dem Argumentintegral
\eqref{buch:analytisch:argument:eqn:argumentintegral0}
und dem früher eingeführten Begriff der Ordnung einer Nullstelle
oder eines Pols hergestellt werden.
Zu diesem Zweck wird zunächst die Funktion $f(z)=z^n$ 
betrachtet und später auf eine beliebige Nullstelle oder einen
Pol einer Funktion verallgemeinern.

%
% Geometrische Interpretation
%
\subsubsection{Geometrische Interpretation}
Das Argumentintegral
\eqref{buch:analytisch:argument:eqn:argumentintegral0}
berechnet die Änderung des Argumentes von $f(z)$, wenn man
einer geschlossenen Kurve folgt.
Für eine Funktion $f(z)$, die in $z=0$ weder einen Nullpunkt noch
einen Pol hat, und eine geschlossene Kurve $\gamma$, die den
Nullpunkte einmal in positiver Richtung einmal umläuft, ist
$f(\gamma(t))$ eine Kurve, entlang der sich das Argument um ein ganzzahliges
Vielfaches von $2\pi$ ändert.
Die ganze Zahl $k$, die in der Formel
\eqref{buch:analytisch:argument:eqn:argumentintegral} auftritt
ist also die Windungszahl des Weges $f(\gamma(t))$.

%
% Berechnung des Integrals
%
\subsubsection{Berechnung des Integrals für $f(z)=z^n$}
Wir betrachten zunächst die Funktionen $f(z)=z^n$ und berechnen 
das Integral
\begin{equation}
\frac{1}{2\pi i}
\oint_{\gamma}
\frac{f'(z)}{f(z)}
\,dz
=
\frac{1}{2\pi i}
\oint_{\gamma}
\frac{nz^{n-1}}{z^n}
\,dz
=
n\cdot
\frac{1}{2\pi i}
\oint_{\gamma}
\frac{1}{z}
\,dz
\label{buch:analytisch:argument:eqn:argumentintegral2}
\end{equation}
für eine Kurve $\gamma$, die den Nullpunkt nicht trifft und ihn
genau einmal im Gegenuhrzeigersinn umläuft.
Da $f'(z)/f(z)$ ausserhalb des Nullpunkts holomorph ist, können
wir den Weg auch durch den Einheitskreis $t\mapsto e^{2\pi it}$
ersetzen und werden den gleichen Wert für das Integral erhalten,
nämlich
\begin{align*}
\frac{1}{2\pi i}
\oint_{\gamma}
\frac{f'(z)}{f(z)}
\,dz
&=
n\cdot
\frac{1}{2\pi i}
\int_0^1
\frac{1}{e^{2\pi it}}
2\pi ie^{2\pi it}
\,dt
=
n
\int_0^1\,dt
=
n.
\end{align*}
Das Argumentintegral
\eqref{buch:analytisch:argument:eqn:argumentintegral2}
liefert also die Ordnung eines Poles oder einer Nullstelle.

%
% Argumentintegral für eine Funktion mit f(0)=0
%
\subsubsection{Argumentintegral für eine Funktion mit $f(0)=0$}
Die letzte Beobachtung soll noch etwas vertieft werden.
Wir betrachten eine holomorphe Funktion $f(z)$ mit einem
singulären Punkt endlicher Ordnung in $0$.
Für eine Nullstelle der Ordnung $n$ kann die Funktion $f$
als
\[
f(z)
=
z^n\, g(z)
\]
geschrieben, werden, wobei $g(z)$ eine holomorphe Funktion
ist, die im Nullpunkt nicht verschwindet.
Da Nullstellen isoliert sind, verschwindet $g(z)$ auch in
einer genügend kleinen Umgebung von $0$ nicht.

Die Ableitung
\[
f'(z)
=
nz^{n-1}g(z) + z^ng'(z).
\]
ist eine holomorphe Funktion.
Das Argumentintegral über einen kleine Kreis $\gamma$ um den
Nullpunkt ist
\begin{align*}
\frac{1}{2\pi i}
\oint_\gamma
\frac{f'(z)}{f(z)}
\,dz
&=
\frac{1}{2\pi i}
\oint_\gamma
\frac{nz^{n-1}g(z)+z^ng'(z)}{z^ng(z)}
\,dz
\\
&=
\frac{1}{2\pi i}
\oint_\gamma
\frac{n}{z}
+
\frac{g'(z)}{g(z)}
\,dz.
\intertext{Der zweite Term im Integral ist eine holomorphe Funktion
auch im Punkt $0$, der Weg $\gamma$ ist also in im Definitionsgebiet
von $g'(z)/g(z)$ zusammenziehbar und das Integral verschwindet.
Es bleibt daher}
&=
n\cdot
\frac{1}{2\pi i}
\oint_\gamma
\frac{1}{z}
\,dz
=
n,
\end{align*}
wie früher bereits ausgerechnet wurde.

%
% Argumentintegral für eine Polstelle
%
\subsubsection{Argumentintegral für eine Polstelle}
Hat die Funktion $f(z)$ einen Pol der Ordnung $n$, dann kann sie als
\[
f(z) = \frac{1}{z^n}g(z)
\]
geschrieben werden, wobei $g$ eine in einer Umgebung von $0$ holomorphe
Funktion ist, die keine Nullstellen oder Pole nahe bei $0$ hat.
Die Ableitung von $f$ ist
\[
f'(z)
=
-
\frac{n}{z^{n+1}}g(z)
+
\frac{1}{z^n}g'(z).
\]
Für diese Funktion ist das Argumentintegral
\begin{align*}
\frac{1}{2\pi i}
\oint_\gamma
\frac{f'(z)}{f(z)}
\,dz
&=
\frac{1}{2\pi i}
\oint_\gamma
\frac{
-
nz^{-(n+1)}g(z)
+
z^{-n}g'(z)
}{
z^{-n}\,g(z)
}
\,dz
\\
&=
\frac{1}{2\pi i}
\oint_\gamma
-
\frac{n}{z}
+
\frac{g'(z)}{g(z)}
\,dz.
\intertext{Wie im Fall einer Nullstelle ist der Weg $\gamma$
im Definitionsbereich des zweiten Terms Integranden zusammenziehbar,
der zweite Term trägt daher nichts zum Integral bei.
Für den ersten Term findet man}
&=
-n\cdot
\frac{1}{2\pi i}
\oint_\gamma
\frac{1}{z}
\,dz.
\end{align*}

%
% Argumentintegral und Ordnung
%
\subsubsection{Argumentintegral und Ordnung}
Die Berechnungen der vorangegangenen Abschnitte zeigen, dass
der Wert des Argumentintegrals für Nullstellen und Pole 
immer durch die Ordnung
\[
\omega(0;f)
=
\oint_{\gamma}
\frac{f'(z)}{f(z)}
\,dz
\]
ausgedrückt werden kann.
Wir haben damit den folgenden Satz bewiesen.

\begin{satz}[Ordnung]
\label{buch:analytisch:argument:satz:ordnung}
Sei $f$ eine holomorphe Funktion mit einer isolierten Singularität
an der Stelle $a\in\mathbb{C}$.
\index{Singularitat@Singularität}%
Sei $\gamma$ ein kleiner Kreis, der den Punkt $z_0$ einmal im
Gegenuhrzeigersinn umläuft.
Er sei zudem so klein, dass er keine anderen Null- oder Polstellen
von $f$ als $z_0$ enthält.
Dann gilt
\[
\omega(z_0;f)
=
\frac{1}{2\pi i}
\oint_{\gamma}
\frac{f'(z)}{f(z)}
\,dz.
\]
\end{satz}
